\section{Architecture}\label{sec:architecture}

\subsection{Components}

Wysp runs entirely on the user's machine. One Python process, the Core, serves the Panel and accepts adapter connections. Each adapter lives inside its host application and dials out to the Core. Figure~\ref{fig:system} shows the parts; this table says what each one is.

\begin{xltabular}{\textwidth}{@{}L{3.0cm}L{2.5cm}Y@{}}
\toprule
\textbf{Component} & \textbf{Language} & \textbf{Role} \\
\midrule
Contract & Python (Pydantic) & Models of every message, JSON Schema export, predicate language, state digest. Single source of truth for the SDKs and the test kit. \\
Wysp Core & Python 3.13 & Gateway, session, orchestrator, board, narration, trust, verification, pedagogy, LLM providers, voice, MCP endpoint. FastAPI and uvicorn. \\
Panel & TypeScript, React 19 & A minimal shell: chat or diagram, captions, voice, permission dialog, lesson and exercise cards. Board drawn with React Flow and ELK layouts, state in Zustand. \\
Adapter SDKs & Python, TypeScript & WebSocket client, action registry with parameter validation, origin tagging, main-thread dispatcher (Python). \\
Contract Test Kit & Python & Plays the Core's part against any adapter and checks the contract. Includes a mock Kanban adapter. \\
Tutors & Python, Kotlin, TypeScript & Blender extension, IntelliJ plugin, Chrome extension, each with its profile. \\
\bottomrule
\end{xltabular}

\subsection{Trust boundaries}

Wysp has one rule for crossing a boundary: data from the other side is data, never an instruction. Figure~\ref{fig:boundaries} shows the zones and what protects each channel.

\begin{figure}[H]
\centering
\begin{tikzpicture}[fig]
  \node[zone=inkpurple, minimum width=15.6cm, minimum height=5.6cm] (local) at (0,0.2) {};
  \node[font=\sffamily\bfseries\scriptsize, text=inkpurple, anchor=north west] at (-7.7,2.85) {The user's machine (Core bound to 127.0.0.1)};
  \node[zone=lilacdeep, minimum width=3.3cm, minimum height=2.6cm] (host) at (-5.6,0.3) {};
  \node[mod, minimum width=2.7cm, minimum height=0.9cm] (app) at (-5.6,0.95) {Host application\\\tiny Blender, IDE, browser};
  \node[mod, minimum width=2.7cm, minimum height=0.9cm] (adp) at (-5.6,-0.3) {Adapter\\\tiny closed action catalog};
  \node[mod, draw=inkpurple, fill=inkpurple!8, minimum width=3.2cm, minimum height=3.3cm, align=center] (core) at (-0.6,0.2) {\textbf{Wysp Core}\\[3pt]\tiny permission gate\\\tiny provenance validator\\\tiny verification\\\tiny audit log\\\tiny redaction};
  \node[zone=pinkdeep, minimum width=3.0cm, minimum height=1.4cm] (pn) at (4.1,1.3) {};
  \node[font=\sffamily\scriptsize] at (4.1,1.3) {Panel\\\tiny browser, JCEF, Blender UI};
  \node[zone=muted, minimum width=3.0cm, minimum height=1.4cm] (ag) at (4.1,-0.8) {};
  \node[font=\sffamily\scriptsize, align=center] at (4.1,-0.8) {Your AI agent (optional)\\\tiny MCP client, local process};
  \node[zone=muted, minimum width=9.4cm, minimum height=1.5cm, dashed] (cloud) at (-2.4,-4.0) {};
  \node[font=\sffamily\scriptsize, align=center] at (-2.4,-4.0) {\textbf{Remote, only if configured}\\\tiny LLM provider (conversation, observation excerpts, code snippets)\\\tiny Groq Whisper (microphone clips, push-to-talk only)};
  \draw[biarr] (adp.east) -- node[lab, above, align=center] {\code{/ws/adapter}\\token + Origin} (core.west |- adp.east);
  \draw[biarr] (core.east |- pn.west) -- node[lab, above] {\code{/ws/ui} Origin} (pn.west);
  \draw[biarr] (core.east |- ag.west) -- node[lab, above, align=center] {\code{/mcp}\\token + Origin} (ag.west);
  \draw[arr, color=pinkdeep] (core.south) -- ++(0,-1.1) -- node[lab, right, pos=0.55] {HTTPS} (cloud.north -| core.south);
\end{tikzpicture}
\caption{Zones and channels. The adapter, the Panel and the agent are separate clients of the Core. Only the dashed zone leaves the machine, and only when the user configured it.}\label{fig:boundaries}
\end{figure}

\begin{xltabular}{\textwidth}{@{}L{3.2cm}L{3.9cm}Y@{}}
\toprule
\textbf{Channel} & \textbf{Protection} & \textbf{Notes} \\
\midrule
\code{WS /ws/adapter} & Pairing token in \code{wysp.hello}; Origin allowlist & Token is \code{~/.wysp/token}, 32 random bytes, created once. A missing Origin (native clients) is allowed. \\
\code{WS /ws/ui} & Origin allowlist & No token. The first Panel is elected as the one that speaks. \\
\code{POST /api/listen/*} & Origin and token header \code{x-wysp-token} & Turns the microphone on, so it needs the token. \\
\code{POST /api/tts}, \code{/api/ocr} & Origin allowlist & Local voice and local text recognition only. No proxying of URLs. \\
\code{/mcp} & Token (Bearer or \code{X-Wysp-Token}), Origin, loopback host check & Served in external mode, or with \code{--mcp}. \\
\bottomrule
\end{xltabular}

The Origin allowlist contains the Core's own origin, the Panel dev server, and anything listed in \code{WYSP\_ALLOWED\_ORIGINS} (the Chrome extension's origin, for instance). The CLI refuses to bind to a non-loopback address unless \code{--unsafe-bind} is given.

\subsection{The data flow of a turn}\label{sec:turn}

A \emph{turn} starts when the user sends a message and ends when Wysp has finished answering. Figure~\ref{fig:turn} shows the six stages, and the list below explains what the code does in each.

\begin{figure}[H]
\centering
\begin{tikzpicture}[fig, node distance=4mm]
  \tikzset{st/.style={draw=lilacdeep, fill=lilac!18, rounded corners=3pt, minimum width=2.45cm, minimum height=1.25cm, align=center, inner sep=3pt}}
  \node[st] (o) {\textbf{1 Observe}\\\tiny adapter.observe\\\tiny survey, focus};
  \node[st, right=of o] (r) {\textbf{2 Reason}\\\tiny LLM + tools\\\tiny at most 8 rounds};
  \node[st, right=of r] (n) {\textbf{3 Narrate}\\\tiny segments + cues\\\tiny captions, speech};
  \node[st, right=of n] (b) {\textbf{4 Board}\\\tiny projection\\\tiny provenance check};
  \node[st, below=9mm of n, fill=pink!22, draw=pinkdeep] (g) {\textbf{5 Guide or do}\\\tiny propose\_action\\\tiny permission gate};
  \node[st, left=of g, fill=pink!22, draw=pinkdeep] (e) {\textbf{execute}\\\tiny adapter.execute\\\tiny undo token};
  \node[st, left=of e] (v) {\textbf{6 Verify}\\\tiny fresh observation\\\tiny predicates, evidence};
  \draw[arr] (o) -- (r); \draw[arr] (r) -- (n); \draw[arr] (n) -- (b);
  \draw[arr] (r.south) |- ($(r.south)+(0,-0.45)$) -| (g.north west);
  \draw[arr] (g) -- (e); \draw[arr] (e) -- (v);
  \draw[arr, dashed] (v.north) -- ++(0,0.5) -| node[lab, above, pos=0.2] {fail: auto-undo, report} (o.south);
\end{tikzpicture}
\caption{One turn. Lilac stages run in the Core's own logic; pink stages go through the permission gate and the adapter. The verify stage feeds its result back to the conversation.}\label{fig:turn}
\end{figure}

\begin{enumerate}[leftmargin=1.4em]
  \item \textbf{Observe.} When an adapter becomes the active one, the survey runs: the Core fetches the action catalog, the context and one full observation, and keeps a bounded summary (at most 4000 characters in the prompt). For each turn the system prompt also carries a fresh observation of the focus (up to 20 entities). Sensitive entities are counted, never named.
  \item \textbf{Reason.} The orchestrator builds the messages (system prompt, the last 12 history messages, the user text) and calls the language model with the tools that the adapter's capabilities allow. It loops for at most 8 rounds. A transient LLM error is retried once after 1.5 s. If the model writes a tool call as prose, the Core tries to recover it. Plain text from the model is narrated, so models that skip the \code{narrate} tool still work.
  \item \textbf{Narrate.} Speech is cut into segments of at most 300 characters. A segment can carry \emph{cues} (``at word 7, focus this node'') so that the board follows the voice. If the user interrupts, unplayed segments are kept and can be resumed.
  \item \textbf{Board.} \code{show\_entities} does not let the model write nodes. The Core \emph{projects} the observed entities the model selected, and the provenance validator downgrades any \pobs{} node without a matching ref in this session's observations. The Panel receives a patch.
  \item \textbf{Guide or do.} \code{propose\_action} validates the parameters against the action's JSON Schema. In guide mode the Core installs a watch with the action's postconditions, highlights the target in the application and waits for the user. In do mode the Core asks the permission gate, shows a dry-run preview and, if approved, executes.
  \item \textbf{Verify.} The Core observes again, evaluates the postconditions and produces \emph{evidence}. If the postconditions fail after a Wysp action, the Core undoes it and reports the failure.
\end{enumerate}

\begin{wherebox}
\file{core/src/wysp\_core/app.py} (endpoints, Origin allowlist), \file{gateway/server.py} (handshake, heartbeat), \file{orchestrator/turn.py} (turn loop), \file{orchestrator/context.py} (prompt), \file{orchestrator/pipeline.py} (guide, do, verify, undo). Specs: \file{.kiro/specs/core-orchestrator/}.
\end{wherebox}



